\documentclass[12pt,a4paper]{article}

% ── Fonts and encoding (LuaLaTeX) ─────────────────────────────────────────────
\usepackage{fontspec}
\usepackage{unicode-math}
\setmainfont{Linux Libertine O}
\setmathfont{TeX Gyre Pagella Math}
\setmonofont[Scale=0.85]{Everson Mono}

% ── Layout ────────────────────────────────────────────────────────────────────
\usepackage[top=1in, bottom=1in, left=1.1in, right=1.1in, headheight=15pt]{geometry}
\usepackage{microtype}
\usepackage{parskip}
\usepackage[english]{babel}
\setlength{\emergencystretch}{3em}

% ── Headers ───────────────────────────────────────────────────────────────────
\usepackage{fancyhdr}
\pagestyle{fancy}
\fancyhf{}
\fancyhead[L]{\leftmark}
\fancyhead[R]{\thepage}
\renewcommand{\headrulewidth}{0.4pt}

% ── Section styling ───────────────────────────────────────────────────────────
\usepackage{titlesec}
\titleformat{\section}{\Large\bfseries}{\thesection}{1em}{}
\titleformat{\subsection}{\large\bfseries}{\thesubsection}{1em}{}

% ── Lists ─────────────────────────────────────────────────────────────────────
\usepackage{enumitem}
\setlist[itemize]{topsep=0.5ex, itemsep=0.25ex, parsep=0pt}
\setlist[enumerate]{topsep=0.5ex, itemsep=0.25ex, parsep=0pt}

% ── Math and tables ───────────────────────────────────────────────────────────
\usepackage{amsmath,amsthm}
\usepackage{mathtools}
\usepackage{booktabs}
\usepackage{array}
\usepackage{tabularx}

% ── Cross-references ──────────────────────────────────────────────────────────
\usepackage{hyperref}
\usepackage{cleveref}
\hypersetup{colorlinks=true, linkcolor=blue, urlcolor=cyan, citecolor=blue,
  pdfkeywords={SIC-POVM, moduli field, ray class field, Stark units, integer norm sieve, algebraic number theory, Zauner conjecture, PARI/GP, d=2048}}

% ── Theorem environments ──────────────────────────────────────────────────────
\newtheorem{theorem}{Theorem}[section]
\newtheorem{lemma}[theorem]{Lemma}
\newtheorem{proposition}[theorem]{Proposition}
\newtheorem{corollary}[theorem]{Corollary}
\newtheorem{definition}[theorem]{Definition}
\newtheorem{remark}[theorem]{Remark}

% ── Custom notation ───────────────────────────────────────────────────────────
\newcommand{\Fd}{F}
\newcommand{\md}{m_d}
\newcommand{\ord}{\mathcal{O}}
\newcommand{\Norm}{\mathrm{N}}
\newcommand{\Lean}{\textsc{Lean\,4}}

\hypersetup{pdfauthor={Anonymous for review}, pdfcreator={}}
\begin{document}

\title{\textbf{The 2048 Chrysopoeia:\\
An Explicit Construction Program for the $d=2048$ SIC-POVM Moduli}}


\date{7 July 2026}

\maketitle

\begin{abstract}
Companion to the SIC--Stark--12th formalization, this paper lays out an explicit program for the
diagonal moduli $N_k=|\psi_k|^2$ of the $d=2048$ SIC-POVM.  It reports what direct
computation has settled and what remains open.  Five results are in place.  The
$a=0$ stratum reduces to a flat autocorrelation condition on the real moduli, a
condition verified on $d=12$ to a residual near $10^{-58}$.  Unconstrained
numerical recovery is blocked by a genuine local minimum (residual
$3.9\times 10^{-3}$).  The moduli field is the real index-$2$ subfield, of
degree $2^{27}$, of the ray class field of conductor $(d)\,\infty_1\infty_2$,
calibrated on the $d=12$ fiducial where the moduli have degree $4$ or $8$ and
generate a field of degree $16$; its unramified ascent is verified to degree
$128$.  Stark-unit $L$-values
supply the moduli without building the full field, a method validated end to end
on $d=12$.  An integer norm sieve separates degenerate $S$-unit aliases by
discriminant $-\md$.  Still open: the reduced character set for the $1024$
moduli and unconditional existence.
\end{abstract}

\noindent\textbf{Keywords:} SIC-POVM \(\cdot\) moduli field \(\cdot\) ray class
field \(\cdot\) Stark units \(\cdot\) integer norm sieve \(\cdot\) algebraic
number theory \(\cdot\) Zauner conjecture \(\cdot\) PARI/GP \(\cdot\) $d=2048$

\newpage
\tableofcontents
\newpage

%──────────────────────────────────────────────────────────────────────────────
\section{Introduction}
\label{sec:intro}
%──────────────────────────────────────────────────────────────────────────────

A symmetric informationally complete positive operator-valued measure
(SIC-POVM) in dimension $d$ is a set of $d^2$ unit vectors with constant
pairwise overlap $|\langle \psi_j, \psi_k\rangle|^2 = 1/(d+1)$ for $j\neq k$.
The Weyl--Heisenberg (WH) construction seeks a single fiducial vector
$\psi\in\mathbb{C}^d$ whose orbit under the $d^2$ displacement operators
$D_{a,b}$ is such a set.  Existence is Zauner's conjecture \cite{Zauner1999},
still open in general and verified numerically through $d\approx 2400$.

Appleby, Bengtsson, Flammia, and collaborators
\cite{Appleby2013,Appleby2017,AFMY2020} showed that the fiducial coordinates
are algebraic numbers of a striking arithmetic: they lie in specific ray class
fields of the real quadratic field $\Fd = \mathbb{Q}(\sqrt{(d+1)(d-3)})$.  This
ties SIC existence directly to the real-quadratic case of Hilbert's twelfth
problem.  The companion paper \cite{StarkHilbert12} formalizes the chain in
\Lean{} and, at the frontier dimension $d=2048=2^{11}$, proves the transport
apparatus sorry-free while leaving unconditional existence as the single open
node.

Here we take the constructive counterpart.  Instead of transporting a fiducial
whose existence is assumed, we construct the moduli $N_k = |\psi_k|^2$ directly
as exact algebraic numbers and report a program whose components are all in
place.  We restrict ourselves to the diagonal $a=0$ stratum, which we show is
governed entirely by the real moduli; the phase strata are deferred to sequel
work.

The account is grounded in computation.  Each field is built in PARI/GP
\cite{PARI}, each identity is exhibited numerically to stated precision, and
the one bare-metal computation, the norm sieve, is executed and its output
reproduced.  We are equally deliberate about the boundary of what is proved: the
final map from $L$-values to the $1024$ moduli is not yet closed, and we say so
plainly in \Cref{sec:status}.  The title names the aim, the making of the gold,
not a completed casting.

%──────────────────────────────────────────────────────────────────────────────
\section{The base field and its $S$-unit generators}
\label{sec:base}
%──────────────────────────────────────────────────────────────────────────────

Fix $d=2048$.  The SIC discriminant is
\[
\md = (d+1)(d-3) = 2049\cdot 2045 = 4{,}190{,}205 = 3\cdot 5\cdot 409\cdot 683,
\]
squarefree, so $\Fd=\mathbb{Q}(\sqrt{\md})$ is real quadratic with
$\operatorname{disc}\Fd = \md$ (as $\md\equiv 1 \bmod 4$).  PARI/GP gives class
number $h(\Fd)=64$, class group $\mathbb{Z}/32\times\mathbb{Z}/2$, and
regulator $R \approx 7.6241$.

The four primes dividing $\md$ ramify in $\Fd$, and the $S$-unit group for
$S = \{3,5,409,683\}$ has rank five.  Its generators, computed explicitly, are
\[
\varepsilon = \tfrac{2047 - \sqrt{\md}}{2},\quad
3,\quad 5,\quad
g_3 = \tfrac{\sqrt{\md}-2045}{2},\quad
g_4 = \tfrac{2049-\sqrt{\md}}{2},
\]
together with the torsion unit $-1$.  Their field norms are
\[
\begin{gathered}
\Norm(\varepsilon)=+1,\quad
\Norm(3)=9,\quad
\Norm(5)=25,\\
\Norm(g_3)=-(d-3)=-2045,\quad
\Norm(g_4)=(d+1)=2049 .
\end{gathered}
\]

Three arithmetic coincidences organise everything that follows.

The SIC denominator is a generator norm: $\Norm(g_4)=d+1$, so $1/(d+1)$, the
value that defines equiangularity, is simply the reciprocal of the field norm
of $g_4$.

The fundamental unit is the mean modulus: at the principal real embedding
$\sqrt{\md}\approx 2046.99902$ one has
$|\varepsilon| = 0.000488519\ldots \approx 1/d = 0.000488281\ldots$, so
$\log|\varepsilon| = -R$ and the coarse scale of every modulus is a single power
of $\varepsilon$.

The discriminant factors through the dimension: $\md = (d+1)(d-3)$, with
$d+1=3\cdot 683$ and $d-3 = 5\cdot 409$, so the two signed prime discriminants
that carry the class-group structure are exactly the factors of $d\pm$ small.

Every modulus, if it is an $S$-unit of $\Fd$, is a monomial
$\pm\,\varepsilon^{a}3^{b}5^{c}g_3^{e}g_4^{f}$ with integer exponents, and its
norm is the integer $9^{b}25^{c}(-2045)^{e}(2049)^{f}$.  This integrality is the
lever of \Cref{sec:sieve}.

%──────────────────────────────────────────────────────────────────────────────
\section{The $a=0$ stratum is a flat autocorrelation condition}
\label{sec:autocorr}
%──────────────────────────────────────────────────────────────────────────────

Let $\omega = e^{2\pi i/d}$ and write the diagonal overlaps
$\nu_{0,b} = \langle \psi, Z^b \psi\rangle = \sum_k N_k\,\omega^{bk}$, where
$Z$ is the clock operator and $N_k=|\psi_k|^2$.  The overlap vector
$(\nu_{0,b})_b$ is the discrete Fourier transform of the moduli vector
$(N_k)_k$.  The SIC condition on the $a=0$ stratum is
$|\nu_{0,b}|^2 = 1/(d+1)$ for all $b\neq 0$, with $\nu_{0,0}=\sum_k N_k = 1$.

\begin{proposition}[Flat autocorrelation]
\label{prop:autocorr}
A real nonnegative vector $(N_k)$ with $\sum_k N_k = 1$ satisfies
$|\nu_{0,b}|^2 = 1/(d+1)$ for every $b\neq 0$ if and only if its cyclic
autocorrelation $C_m = \sum_k N_k N_{k+m}$ obeys
\[
C_0 = \sum_k N_k^2 = \frac{2}{d+1},
\qquad
C_m = \frac{1}{d+1}\ \ (m\neq 0).
\]
\end{proposition}

\begin{proof}
By the Wiener--Khinchin identity $|\nu_{0,b}|^2 = \sum_m C_m\,\omega^{bm}$, the
sequence $(|\nu_{0,b}|^2)_b$ is the transform of $(C_m)_m$.  Imposing
$|\nu_{0,0}|^2 = 1$ and $|\nu_{0,b}|^2 = 1/(d+1)$ for $b\neq 0$ and inverting the
transform gives $C_0 = \tfrac1d\big(1 + \tfrac{d-1}{d+1}\big) = \tfrac{2}{d+1}$
and $C_m = \tfrac1d\big(1 - \tfrac{1}{d+1}\big) = \tfrac1{d+1}$ for $m\neq 0$.
\end{proof}

The diagonal stratum therefore involves no phases and no complex field; it is a
system of real quadratic equations on the moduli.  Two structural facts reduce it
further.  The moduli come in Galois-conjugate pairs under the two real
embeddings of $\Fd$,
\begin{equation}
\label{eq:galois-pair}
N_{k+d/2} = \sigma(N_k),
\end{equation}
so the $2048$ moduli are determined by $1024$ field elements and their
conjugates.  The whole reformulation is exact.

\begin{remark}[Validation on $d=12$]
\label{rem:d12-autocorr}
The known $d=12$ fiducial has moduli that are exact elements of a real number
field \cite{Appleby2013}.  Evaluating them at the chosen embedding and forming
the autocorrelation, one finds $\sum_k N_k = 1$, $C_0 = 2/13$, and
$C_m = 1/13$ for all $m\in\{1,\dots,11\}$, each to a residual below $10^{-58}$
at working precision; the pairing \eqref{eq:galois-pair} with $d/2=6$ holds to
the same precision.  \Cref{prop:autocorr} is thus not only proved but confirmed
against ground truth.
\end{remark}

%──────────────────────────────────────────────────────────────────────────────
\section{The numerical assay fails structurally}
\label{sec:numeric}
%──────────────────────────────────────────────────────────────────────────────

We record why numerical search cannot produce the moduli; it is the natural
first attempt, and it has stalled the frontier.  We formed the full SIC objective
$\Phi(\psi) = \sum_{p\neq 0}\big(|\langle\psi,D_p\psi\rangle|^2 - 1/(d+1)\big)^2$
over all $d^2-1 = 4{,}194{,}303$ overlaps, evaluated by fast Fourier transform,
and minimized it from the best available numerical seeds, including seeds
confined to the Zauner-symmetric subspace where high-precision fiducials are
normally found.

\begin{proposition}[Spurious minimum]
\label{prop:spurious}
The best Zauner-symmetric seed is a genuine local minimum of $\Phi$ at residual
$\max_p\big||\langle\psi,D_p\psi\rangle|^2 - 1/(d+1)\big| = 3.87\times 10^{-3}$:
its gradient norm is $1.96\times 10^{-8}$ and every line-search direction
increases $\Phi$.
\end{proposition}

No first- or second-order local method escapes a genuine local minimum, so the
seeds are exhausted, not merely unpolished.  A second obstruction is precision:
recognizing a modulus as a specific $S$-unit monomial by integer-relation
detection against the generator logarithms requires the modulus to about twelve
significant digits, whereas a fiducial at SIC residual $3.9\times 10^{-3}$ pins
its moduli to two or three.  The construction must therefore be algebraic, and
the moduli must be produced as exact special values rather than recovered from
floating-point fits.

%──────────────────────────────────────────────────────────────────────────────
\section{The moduli field: a pro-$2$ ray class tower}
\label{sec:tower}
%──────────────────────────────────────────────────────────────────────────────

\subsection{Sizing, calibrated on $d=12$}

The moduli are real at the fiducial embedding, but reality at one embedding does
not force total reality, and the $d=12$ ground truth shows it fails.  The correct
field is the index-$2$ real subfield of the ray class field of $\Fd$ at conductor
$(d)\,\infty_1\infty_2$ (both archimedean places ramified), fixed by complex
conjugation at the fiducial embedding.  We fix the rule by calibration against
$d=12$, where the moduli are known exactly from the machine-checked existence
ring of the companion paper \cite{StarkHilbert12}.

\begin{proposition}[$d=12$ moduli degrees and field]
\label{prop:d12-deg}
The moduli of the $d=12$ fiducial have minimal polynomials of degree $4$ or $8$
over $\mathbb{Q}$: degree $4$ for $N_0, N_3, N_6, N_9$, with
\[
97344\,x^4 - 32448\,x^3 + 3224\,x^2 - 104\,x + 1 = 0
\]
the minimal polynomial of $N_0$ and $N_6$, and degree $8$ for the remaining
eight.  The field generated by all twelve moduli has degree $16$ over
$\mathbb{Q}$: it is the real index-$2$ subfield of the
conductor-$(12)\,\infty_1\infty_2$ ray class field of $\mathbb{Q}(\sqrt{13})$
(degree $32$), of signature $(8,4)$, hence not totally real.
\end{proposition}

Each claim is verified in PARI/GP directly against the exact power-basis moduli
vectors of the companion formalization.  This corrects two naive expectations at
once.  First, the moduli do not sit in a proper subfield of the degree-sixteen
real part of the coordinate field: individual moduli have degree only $4$ or
$8$, but together they generate the full degree-$16$ field.  Second, real is
not totally real: the totally-real ray class field of conductor $(12)$ (degree
$8$ over $\mathbb{Q}$, ray class group $(\mathbb{Z}/2)^2$) is a proper subfield
of the moduli field, and the archimedean ramification carries the rest.  The
corrected rule: the moduli field is the real index-$2$ subfield of the ray
class field of conductor $(d)\,\infty_1\infty_2$, so its degree over
$\mathbb{Q}$ equals the ray class order $|\mathrm{Cl}_{(d)\infty_1\infty_2}|$;
for $d=12$ that order is $16$, in exact agreement with the machine-checked
moduli.  Apply it to $d=2048$ and PARI/GP gives ray class group
\[
[\,4096,\ 512,\ 8,\ 4,\ 2\,],\qquad \text{order } 2^{27},
\qquad \text{degree } 2^{27}\ \text{over }\mathbb{Q}
\]
(the conductor-$(2048)$ group without the archimedean places,
$[\,4096,\ 512,\ 8,\ 2\,]$ of order $2^{25}$, sizes the totally-real
sub-tower).  Every cyclic factor is a power of two: the moduli field is a pure
pro-$2$ tower, matching $d=2^{11}$.  It is a large field, and its defining
polynomial, with order $2^{27}$ coefficients, is not a writable object.  This
is the central constraint the remainder of the paper is built around.

\subsection{The unramified ascent, verified layer by layer}

The monolithic route is closed at once: the Stark-unit computation of the full
Hilbert class field, \texttt{quadhilbert} in PARI/GP, exhausts memory beyond
eight gigabytes for $h=64$.  The tower must be climbed one quadratic step at a
time, and the unramified part, of degree $h(\Fd)=64$ over $\Fd$, can be.  We
verify the following ascent, each level by \texttt{bnrclassfield} on the Hilbert
ray class group.

\begin{center}
\begin{tabular}{@{}clll@{}}
\toprule
Level & Field & $\deg/\mathbb{Q}$ & Construction \\
\midrule
$0$ & $\Fd = \mathbb{Q}(\sqrt{\md})$ & $2$ & base, $h=64$, class $[32,2]$ \\
$1$--$2$ & genus $\mathbb{Q}(\sqrt5,\sqrt{409},\sqrt{2049})$ & $8$ & $(\mathbb{Z}/2)^2$, unramified \\
$3$ & $C_4$ & $16$ & R\'edei $\md = 409\cdot 10245$ \\
$4$ & $C_8$ & $32$ & \texttt{bnrclassfield}, contains $C_4$ \\
$5$ & $C_{16}$ & $64$ & \texttt{bnrclassfield} \\
$6$ & $C_{32}$, Hilbert class field & $128$ & \texttt{bnrclassfield}, $h$ reached \\
\bottomrule
\end{tabular}
\end{center}

The genus field is generated by square roots of the $S$-unit generator norms:
$\sqrt{2049} = \sqrt{d+1}$ and $\sqrt{2045}=\sqrt{d-3}=\sqrt{5}\sqrt{409}$.  Its
discriminant is $\md^{4}$, so the relative discriminant over $\Fd$ is trivial and
the extension is unramified.  The first non-genus step is fixed by R\'edei theory
on the signed prime discriminants $\{-3,5,409,-683\}$: the $\mathbb{F}_2$
R\'edei matrix has rank two, giving $4$-rank one in agreement with the class
group $[32,2]$, and the associated factorization $\md = 409\times 10245$ (with
$10245 = 3\cdot 5\cdot 683$) names the cyclic quartic.  The prime it isolates,
$409$, divides $d-3$.  The ascent reaches the full Hilbert class field of degree
$128$ over $\mathbb{Q}$; the ramified part of the tower, of index $2^{20}$ over
the Hilbert class field, remains to be climbed by the same method, and we do not
require it explicitly for what follows, because of the value route of
\Cref{sec:portals}.

%──────────────────────────────────────────────────────────────────────────────
\section{The moduli as $L$-function special values}
\label{sec:portals}
%──────────────────────────────────────────────────────────────────────────────

Since the degree-$2^{27}$ moduli field cannot be written as a polynomial, the
moduli must be obtained as numbers, not as coordinates in a field basis.  The
mixed-signature Stark conjecture supplies them: for each character $\chi$ of the
ray class group, the derivative $L'(0,\chi)$ of the associated Hecke
$L$-function is the logarithm of a Stark unit, an algebraic number in the class
field, and the moduli are assembled from these units.  The functional equation
relates $s=0$, where the algebraic unit lives, to $s=1$, where the transcendental
analytic value lives; one reaches the algebraic point at $s=0$ directly.

Two remarks make this operational.  First, a real modulus is a sum over a
conjugate pair of characters $\chi,\bar\chi$, so
$L'(0,\chi)+L'(0,\bar\chi)$ is real by construction: the imaginary parts cancel.
Second, one never forms the field.  The special values are computed to arbitrary
precision from the $L$-function directly.

\begin{proposition}[Value route, validated on $d=12$]
\label{prop:portals-d12}
For $\mathbb{Q}(\sqrt{13})$ at the $d=12$ conductor, \texttt{bnrstark} returns
the totally-real sub-tower at conductor $(12)$ as a degree-four relative
polynomial over $\mathbb{Q}(\sqrt{13})$ generated by the Stark unit, and
\texttt{bnrL1} returns the derivatives $L'(0,\chi)$ for all sixteen characters
of the conductor-$(12)\,\infty_1\infty_2$ group to sixty digits.  The moduli
field of \Cref{prop:d12-deg} is thereby reachable from $L$-values without any
polynomial of prohibitive degree: the even characters deliver its totally-real
sub-tower, and the odd (mixed-signature) characters, the ones that see the
archimedean ramification, carry the remaining index: exactly the regime of
the mixed-signature Stark conjecture this program assumes.
\end{proposition}

The same machinery applies at $d=2048$; the only obstruction is scale.  An
exhaustive computation over all $2^{25}$ characters is infeasible, but the $a=0$
moduli are only $1024$ real numbers, pinned by the flat autocorrelation of
\Cref{prop:autocorr} and paired by \eqref{eq:galois-pair}.  The open task is the
determination of the reduced set of character orbits whose $s=0$ derivatives
carry the diagonal stratum; we return to it in \Cref{sec:status}.

%──────────────────────────────────────────────────────────────────────────────
\section{The integer norm sieve}
\label{sec:sieve}
%──────────────────────────────────────────────────────────────────────────────

A special value delivers a number to high precision.  A number is not yet an
identity: at any finite precision, a real value may fit more than one $S$-unit
monomial.  We resolve this last ambiguity exactly, using the integrality
recorded at the end of \Cref{sec:base}.

The two fine generators $g_3,g_4$ have principal-embedding magnitudes on either
side of $1$, with nearly opposite logarithms,
\[
\begin{gathered}
\log|g_3| = -0.000488639\ldots,\qquad
\log|g_4| = +0.000488401\ldots,\\
\log|g_3|+\log|g_4| = -2.387\times 10^{-7}.
\end{gathered}
\]
Their product is therefore a magnitude near-unit.  Two monomials that differ by
the factor $g_3 g_4$, that is by incrementing both fine exponents, are
magnitude-identical to about one part in $10^{7}$: this is the ambiguity that
numerical fitting cannot break.  But their field norms differ by exactly
\[
\Norm(g_3 g_4) = -(d-3)(d+1) = -\md = -4{,}190{,}205 .
\]

\begin{proposition}[Norm separation]
\label{prop:sieve}
Let $u = \varepsilon^{a}3^{b}5^{c}g_3^{e}g_4^{f}$ and
$u' = \varepsilon^{a}3^{b}5^{c}g_3^{e+1}g_4^{f+1}$.  Then
$\big||u|-|u'|\big|/|u| < 3\times 10^{-7}$ at the principal embedding, while
$\Norm(u')/\Norm(u) = -\md$.  The exact integer norm separates $u$ from $u'$
though their magnitudes do not.
\end{proposition}

The construction is thus over-determined by three independent constraints, and
their intersection is a unique identity: the portal magnitude to about twelve
digits, the exact integer field norm, and the flat autocorrelation holding
jointly across all $1024$ moduli.  Fitting alone is degenerate; the norm, an
integer computed without any floating point, removes the degeneracy; the
autocorrelation ties the $1024$ choices together.

We implemented the norm sieve directly, computing the generators at the
principal embedding and comparing the canonical monomial
$u=\varepsilon^{1}$ (norm $1$, magnitude $0.000488519901637$) with its
magnitude alias $u'=\varepsilon g_3 g_4$ (norm $-4{,}190{,}205$, magnitude
$0.000488519785051$).  The relative magnitude gap is $2.387\times 10^{-7}$, below
recognition tolerance, while the integer norms differ by the full discriminant.
The degeneracy is broken exactly, in integer arithmetic, on the machine.

%──────────────────────────────────────────────────────────────────────────────
\section{Status and open problems}
\label{sec:status}
%──────────────────────────────────────────────────────────────────────────────

We separate what is established from what is not, in the manner of the companion
paper's honest ledger.

Computation has established several facts.  The $a=0$ stratum is the flat
autocorrelation condition of \Cref{prop:autocorr}, verified against $d=12$ to a
residual near $10^{-58}$ (\Cref{rem:d12-autocorr}).  The numerical route is
structurally blocked by a genuine local minimum (\Cref{prop:spurious}).  The
moduli field is the real index-$2$ subfield of the
conductor-$(d)\,\infty_1\infty_2$ ray class field, of degree $2^{27}$ over
$\mathbb{Q}$, calibrated on the exact $d=12$ moduli (degrees $4$ and $8$,
moduli field degree $16$; \Cref{prop:d12-deg}); its unramified part is verified layer by layer to the
Hilbert class field of degree $128$ (\Cref{sec:tower}).  The value route via
Stark units and $L'(0,\chi)$ produces the moduli field without a prohibitive
polynomial, validated end to end on $d=12$ (\Cref{prop:portals-d12}).  The
recognition degeneracy is resolved by the exact norm sieve (\Cref{prop:sieve}),
implemented and executed.

Three problems remain open.  First, the reduced set of ray class characters
whose $s=0$ derivatives determine the $1024$ independent $a=0$ moduli: this is
the one computation that stands between the present apparatus and the explicit
moduli.  Second, the ramified ascent of the tower above the Hilbert class field
(index $2^{20}$) may become necessary if the reduced character route requires
the full conductor.  Third, consequently, unconditional existence
$\texttt{SICPOVM\_Exists}\,2048$, which is the open Zauner conjecture of
\cite{StarkHilbert12}, and the phase strata $a\neq 0$ beyond the diagonal,
remain open.

The program is closed in its parts and validated on the dimension where the
answer is known.  The gold is named, covered in 2048 leaves, and the furnace is
built; the reduced portal computation is the last distillation, and it is a
computation, not a wall.

%──────────────────────────────────────────────────────────────────────────────
\section*{Artifact statement}
The companion formalization \cite{StarkHilbert12} and its Lean modules live in
the \texttt{p4ramill} library of \texttt{p4rakernel}
(\url{https://github.com/umpolungfish/p4rakernel}), frozen at commit
\texttt{eea2c0c} on branch \texttt{crystalline/manuscripts3-2026-07-07} (tag
\texttt{crystalline-manuscripts3-v1}).  This manuscript trio is frozen in
\texttt{ig-docs} (\url{https://github.com/umpolungfish/ig-docs}) on the same
branch name and tag; \texttt{main} continues in each repository.

%──────────────────────────────────────────────────────────────────────────────
\section*{Acknowledgements}
%──────────────────────────────────────────────────────────────────────────────
The author would like to thank Harry T.\ Larson for imparting the importance
of catching rising problems and never letting them go \cite{Larson1961}.
The $d=2048$ moduli program was undertaken in the same spirit: when
unconstrained recovery stalled at a spurious local minimum, the route through
Stark units, the ray class tower, and the integer norm sieve was not dropped.
The field computations were performed in PARI/GP \cite{PARI}; the companion
formalization is \cite{StarkHilbert12}.

%──────────────────────────────────────────────────────────────────────────────
\begin{thebibliography}{9}
%──────────────────────────────────────────────────────────────────────────────

\bibitem{Zauner1999}
G.\ Zauner, \emph{Quantendesigns: Grundz\"uge einer nichtkommutativen
Designtheorie}, Ph.D.\ thesis, Universit\"at Wien, 1999.

\bibitem{Appleby2013}
D.\ M.\ Appleby, H.\ Yadsan-Appleby, and G.\ Zauner,
\emph{Galois automorphisms of a symmetric measurement},
Quantum Inf.\ Comput.\ \textbf{13} (2013), 672--720.

\bibitem{Appleby2017}
M.\ Appleby, S.\ Flammia, G.\ McConnell, and J.\ Yard,
\emph{SICs and algebraic number theory},
Found.\ Phys.\ \textbf{47} (2017), 1042--1059.

\bibitem{AFMY2020}
M.\ Appleby, S.\ Flammia, G.\ McConnell, and J.\ Yard,
\emph{Generating ray class fields of real quadratic fields via complex
equiangular lines}, Acta Arith.\ \textbf{192} (2020), 211--233.

\bibitem{StarkHilbert12}
\emph{SIC-POVMs, a Stark Conjecture, and the 12th:
A Formalization via Paraconsistent Belnap Multilattices}, companion
manuscript, 2026.

\bibitem{PARI}
The PARI~Group, \emph{PARI/GP version 2.13}, Univ.\ Bordeaux, 2021,
\texttt{https://pari.math.u-bordeaux.fr/}.

\bibitem{Larson1961}
H.\ T.\ Larson,
\textit{Catch a Rising Problem\ldots\ and Never Ever Let it Go},
\emph{IRE Transactions on Engineering Management} \textbf{8}(4) (1961), 173--174.
\url{https://ieeexplore.ieee.org/stamp/stamp.jsp?arnumber=1641382}

\end{thebibliography}

\end{document}
